\cleardoublepage

\section{WREC 专家缓存策略设计与实现}

第 3 章将资源受限 MoE 推理建模为专家引用序列上的在线缓存问题。已有 MoE 推理系统通常从专家卸载、异步预取、放置策略或执行流水线角度降低专家权重搬移成本~\cite{xue2024moeinfinity,kamahori2024fiddler,cao2025moelightning,eliseev2023offloading}，也有工作尝试利用模型内部表示预测未来专家访问以支持预取~\cite{zhu2025preattention}。本章在该问题背景下给出 WREC（Workload-aware Router-event Expert Cache）专家缓存策略的设计与实现。WREC 从历史路由轨迹中构造工作负载先验，在在线事件流中维护请求级和词元级状态，并在缓存未命中时基于评分函数执行专家准入、绕过与换出。该策略不修改模型参数、路由器输出和生成结果，仅调整专家权重的驻留集合。

本章介绍 WREC 的方法设计与实现边界。WREC 面向部分驻留 MoE 推理场景，通过历史路由轨迹构造工作负载先验，并在在线执行过程中维护请求级与词元级状态，用于专家缓存的准入、绕过和换出决策。围绕这一流程，本章依次说明设计目标与信息边界、路由事件处理与先验构造、在线评分函数与 WREC-H2 配置，并给出基线策略、轨迹驱动重放模拟器和运行时接口的实现范围。
为区分通用框架与本文主实验配置，后文将评分函数 \eqref{equ:wrec-score} 所定义的方法族称为 WREC，将参数设为 \eqref{equ:wrec-main-config} 的实例称为 WREC-H2。WREC-H2 保留请求内局部项和跨层转移项，关闭最近历史项与驻留竞争惩罚项，是第 5 章主要评估的在线策略。


\begin{figure}[htbp]
    \centering
    \includegraphics[width=0.51\linewidth]{caching_strategy_flowchart_with_scoring_function.png}
    \caption{WREC 专家缓存策略的整体流程}
    \label{fig:wrec-overall-flow}
\end{figure}

\subsection{设计目标与信息边界}

WREC 面向专家传输受限、部分专家驻留的 MoE 预填充阶段场景。在该场景中，GPU 显存无法稳定容纳全部专家对象，系统需要在有限专家缓存容量下选择一部分专家驻留。类似场景在 FlexGen、MoE-Infinity 和 MoE-Lightning 等资源受限推理系统中均被视为影响吞吐和延迟的重要因素~\cite{sheng2023flexgen,xue2024moeinfinity,cao2025moelightning}。当访问专家已经驻留时，推理过程可以直接执行对应 FFN；当访问专家未驻留时，系统需要从 CPU 内存加载专家权重，并可能产生额外传输等待。因此，WREC 的直接目标是在固定专家缓存预算下减少需求加载次数和专家传输量。

WREC 的在线决策仅依赖两类信息：一类是由训练侧或历史工作负载轨迹统计得到的先验，包括专家热度、短窗口出现概率和跨层转移概率；另一类是当前请求已经产生的路由事件，包括当前词元在较早层的专家选择。利用历史轨迹或激活统计指导专家放置，与 MoE-Infinity 中基于专家激活信息减少数据搬移的思想相近~\cite{xue2024moeinfinity}；但 WREC 明确限制在线阶段不可读取评估轨迹中的未来访问。策略重放时，WREC 不读取评估轨迹中的未来访问，也不使用当前词元后续层尚未产生的路由结果。Belady 离线上界和路由窗口 oracle 单独作为离线上界或压力基线处理，不计入可部署在线策略。

WREC 采用准入--绕过机制处理未命中访问。类似的思想在资源受限缓存与卸载系统中较为常见：并非所有被访问对象都值得长期驻留，低复用对象直接进入缓存可能造成缓存污染~\cite{song2023powerinfer,zhang2024h2o}。系统在执行当前 MoE 层计算前需要获得对应专家权重，但该专家不一定进入长期驻留集合。当候选专家评分高于当前缓存中的最低分专家时，WREC 准入该专家并换出低分对象；否则，该专家仅服务当前访问，访问结束后不占用缓存槽。该机制用于降低一次性专家访问对有限缓存的污染。

\subsection{路由轨迹采集与专家引用展开}

WREC 的输入来自 MoE 路由事件。对于支持输出 router logits 的模型，本文在预填充阶段对请求进行前向计算，读取每一层路由器输出，并根据模型配置选择 top-$k$ 个专家。该处理方式与 Mixtral、Switch Transformer 等模型中的稀疏路由机制保持一致~\cite{jiang2024mixtral,fedus2022switch}。一次词元级路由事件表示为

\begin{equation}
    \label{equ:wrec-router-event-detail}
    r=(q,p,l,\mathcal{E}_{q,p,l},\mathcal{P}_{q,p,l}),
\end{equation}

其中，$q$ 表示请求编号，$p$ 表示词元在请求中的位置，$l$ 表示 MoE 层编号，$\mathcal{E}_{q,p,l}$ 表示该词元在该层被选中的专家集合，$\mathcal{P}_{q,p,l}$ 表示对应路由概率。对于采用 top-$2$ 路由的 Mixtral，一个路由事件会展开为两个专家引用。

为了与第 3 章的问题建模保持一致，本文将路由事件展开为按时间排序的专家引用序列。每个引用对象记为 $(l,e)$，其中 $l$ 为层编号，$e$ 为该层中的专家编号。采用 $(l,e)$ 而不是单独的专家编号 $e$ 是必要的，因为不同 MoE 层中的同名专家并不共享权重。例如，第 $3$ 层专家 $1$ 与第 $4$ 层专家 $1$ 是两个不同的缓存对象。

轨迹采集过程保留请求编号、阶段、层号、词元位置、被选专家、路由概率和事件编号等字段。阶段字段用于区分预填充与解码；本文主实验只使用预填充轨迹。事件编号用于验证运行时事件流与离线重放序列之间的顺序一致性。采集后的 JSONL 轨迹被转换为内部 \texttt{ExpertRef} 表示，每条记录包含引用序号、事件序号、请求编号、层号、词元位置和专家编号。

\subsection{工作负载先验构造}

WREC 使用训练侧路由轨迹构造三类工作负载先验。使用工作负载统计作为系统策略输入，是大模型推理调度和资源管理中的常见做法；例如，调度系统可利用长度或历史信息改进请求排序，卸载系统也可利用激活统计指导权重放置~\cite{fu2024efficient,gong2024pastfuture,xue2024moeinfinity}。第一类是短窗口出现概率。设第 $l$ 层按时间顺序排列的路由事件窗口大小为 $w$，专家 $e$ 在一个窗口中至少出现一次的概率定义为

\begin{equation}
    \label{equ:wrec-window-prob}
    P_{\mathrm{win}}(l,e)=
    \frac{\#\{\mathrm{window}:e\in \mathrm{window}\}}{\#\{\mathrm{window}\}}.
\end{equation}

该指标反映专家在短时间范围内被再次访问的可能性。若某专家在多个窗口中频繁出现，则保留该专家可能减少后续需求加载。这一设计对应缓存系统中的时间局部性假设，但对象从普通内存块变为由层号和专家编号共同定义的 MoE 专家对象。

第二类是窗口内期望路由次数。设 $C_{\mathrm{win}}(l,e)$ 表示专家 $e$ 在第 $l$ 层一个窗口中被引用的次数，则

\begin{equation}
    \label{equ:wrec-window-count}
    E_{\mathrm{win}}(l,e)=\mathbb{E}[C_{\mathrm{win}}(l,e)].
\end{equation}

结合单次专家加载传输时间 $T_{\mathrm{load}}$，本文为每个专家对象构造基础收益分

\begin{equation}
    \label{equ:wrec-base-score}
    B(l,e)=P_{\mathrm{win}}(l,e)\cdot E_{\mathrm{win}}(l,e)\cdot T_{\mathrm{load}}-T_{\mathrm{load}}.
\end{equation}

该分数近似表示：若加载并保留一个专家，未来短窗口内可能节省的传输时间与当前加载成本之间的差值。本文主实验中同一模型内各专家大小相同，因此 $T_{\mathrm{load}}$ 在同一配置下为常数；保留该形式是为了使方法能够扩展到专家大小或传输代价不同的场景。

第三类是跨层转移概率。对于同一请求、同一词元，若第 $l-1$ 层选择了专家 $u$，第 $l$ 层选择了专家 $e$，则统计二者在训练轨迹中的共现次数，并定义

\begin{equation}
    \label{equ:wrec-cross-prob}
    P_{\mathrm{cross}}^{(l)}(e\mid u)=
    \frac{N^{(l)}(u,e)}{\sum_{e'}N^{(l)}(u,e')}.
\end{equation}

当某个前层专家在训练轨迹中没有足够样本时，实现中使用均匀分布作为退化处理，以避免未见组合导致评分异常。跨层转移概率的作用是利用当前词元已经产生的上一层路由结果，估计下一层哪些专家更可能被访问。

\subsection{在线状态维护}

在线阶段，WREC 维护三类状态。第一类是每层最近访问历史。系统为每一层维护长度为 $h$ 的专家引用队列，并同步维护队列内专家计数。由此可得到最近访问概率

\begin{equation}
    \label{equ:wrec-recent-prob}
    R_t(l,e)=\frac{\#\{e\ \mathrm{in\ recent\ queue\ of\ layer}\ l\}}{h}.
\end{equation}

该状态主要用于消融实验和扩展配置。第 5 章结果显示，在本文的预填充总槽位预算设置中，最近历史项未形成稳定正贡献，因此主策略将其权重设为 $0$。

第二类是当前请求内统计。对于正在处理的请求，WREC 按层维护该请求已经访问过的专家计数和总引用数。请求内概率定义为

\begin{equation}
    \label{equ:wrec-request-prob}
    Q_t(l,e)=
    \frac{\#\{e\ \mathrm{visited\ in\ current\ request,\ layer}\ l\}}
    {\#\{\mathrm{expert\ refs\ in\ current\ request,\ layer}\ l\}}.
\end{equation}

该项刻画同一输入请求内部的局部复用。与全局静态热度不同，请求内统计会随当前输入内容变化；与 LRU 不同，它在评分时显式区分层号和专家编号。考虑到已有研究指出不同 MoE 模型的局部路由一致性存在差异~\cite{liang2026notallmodels}，该项并不假设所有模型都具有同等稳定的请求内复用，而是在具体轨迹上进行在线统计。

第三类是词元层级历史。WREC 在处理同一词元的不同层时，记录已经完成层的专家选择。当第 $l$ 层的专家 $e$ 被评分时，若同一词元在第 $l-1$ 层选择了专家集合 $\mathcal{E}_{l-1,p}$，则跨层分数为

\begin{equation}
    \label{equ:wrec-cross-score}
    X_t(l,e)=
    \frac{1}{|\mathcal{E}_{l-1,p}|}
    \sum_{u\in \mathcal{E}_{l-1,p}} P_{\mathrm{cross}}^{(l)}(e\mid u).
\end{equation}

该项仅使用当前词元已完成层的路由结果，符合前述在线信息边界。

\subsection{WREC-H2 评分函数与缓存决策}

WREC 对候选专家对象 $(l,e)$ 的评分定义为

\begin{equation}
    \label{equ:wrec-score}
    \mathrm{Score}_t(l,e)=
    B(l,e)+
    \alpha R_t(l,e)+
    \beta Q_t(l,e)+
    \gamma X_t(l,e)-
    \delta D_t(l,e),
\end{equation}

其中，$B(l,e)$ 是训练轨迹构造的基础分，$R_t(l,e)$ 是最近访问项，$Q_t(l,e)$ 是请求内局部项，$X_t(l,e)$ 是跨层转移项，$D_t(l,e)$ 是可选的驻留竞争惩罚项。本文主实验采用的配置为

\begin{equation}
    \label{equ:wrec-main-config}
    w=4,\quad h=8,\quad
    \alpha=0,\quad
    \beta=1024,\quad
    \gamma=1024,\quad
    \delta=0.
\end{equation}

其中，$w$ 表示训练侧窗口大小，$h$ 表示 最近历史长度。根据第 5 章的消融结果，最近历史项在当前预填充总槽位预算重放中未提供稳定收益，因此主配置令 $\alpha=0$，主要依赖请求内统计和跨层转移信号。

当访问 $(l,e)$ 且该专家已经驻留时，WREC 将其记为命中，并更新最近访问、请求内统计和词元层级历史。当访问未命中时，系统首先将该次访问计为一次需求加载，因为真实执行必须获得该专家权重才能完成当前 MoE 层计算；随后 WREC 决定是否让该专家进入长期驻留集合。如果缓存未满，则新专家直接进入驻留集合；如果缓存已满，则 WREC 计算当前驻留专家的分数，选取分数最低者作为候选换出对象。若新专家分数高于该候选换出对象，则执行换出并准入新专家；否则，新专家在本次需求加载后被绕过，不进入驻留集合。

准入--绕过机制与 LRU 的关键区别在于未命中后的驻留决策。LRU 总是在未命中后将新对象放入缓存并换出最久未访问对象；WREC 则允许低分专家在完成当前访问后被绕过。该设计与专家卸载系统中区分短期需求加载和长期驻留放置的思路一致~\cite{xue2024moeinfinity,cao2025moelightning}。该设计适用于 MoE 预填充轨迹中偶发专家访问较多的情形，避免其占用全局有限的专家缓存槽位。

\subsection{基线策略与离线上界}

本文实现五类对照策略。设置多类在线基线、静态基线和离线上界，有助于区分策略本身收益与理想化未来信息带来的上界收益；类似的分层对照也常见于 LLM 服务和资源受限推理系统评估~\cite{kwon2023pagedattention,sheng2023flexgen,cao2025moelightning}。第一，\texttt{on\_demand} 不保留专家缓存，每次专家访问均触发需求加载，用于刻画无缓存时的传输压力。第二，LRU 按在线最近访问顺序换出最久未访问专家。第三，\texttt{static\_hot} 根据训练轨迹中的全局专家频率初始化驻留集合，评估静态热度先验的作用。第四，\texttt{route\_window\_prefetch} 使用评估轨迹中的未来窗口进行预取，仅作为含未来信息的压力基线。第五，Belady 离线上界使用完整未来访问序列换出最晚再次使用或不再使用的专家，作为等大小缓存对象假设下的离线上界。

在容量分配上，本文区分逐层容量和总槽位预算容量。逐层容量为每个 MoE 层分配相同专家槽数，便于与传统分层缓存直觉比较；总槽位预算容量只规定全模型专家槽总数，允许热点层和热点专家在全局范围内竞争槽位。面向批处理 MoE 推理的差异化专家缓存研究也表明，专家驻留决策会受到请求优先级、批次组织和异构内存层次影响，不能简单退化为各层平均分配固定槽位~\cite{li2025diffmoe}。WREC 主结果使用总槽位预算模式，因为该模式更接近显存预算的真实约束：GPU 显存只关心总驻留专家权重大小，而不要求每一层拥有完全相同的槽数。

\subsection{重放模拟器实现}

本文实现了基于访问轨迹的重放模拟器，用于隔离评估缓存策略本身。轨迹驱动重放可以在固定模型输出和固定访问序列下比较不同内存管理策略，避免端到端系统中请求调度、算子实现和运行时负载波动带来的干扰~\cite{sheng2023flexgen,kwon2023pagedattention,xue2024moeinfinity}。模拟器输入为评估路由轨迹、训练路由轨迹、模型专家大小、CPU-GPU 有效带宽和缓存预算。对每条专家引用，模拟器按策略更新缓存状态，并统计需求命中、需求未命中、需求加载、预取加载、预取浪费、传输字节数和传输时间估计值。

对无预取策略而言，需求加载次数等于需求未命中次数。传输时间估计值只计算需求加载字节：

\begin{equation}
    \label{equ:wrec-transfer-time}
    T_{\mathrm{transfer}}=
    \frac{N_{\mathrm{demand}}\cdot B_{\mathrm{expert}}}{b_{\mathrm{cg}}},
\end{equation}

其中，$N_{\mathrm{demand}}$ 是需求加载次数，$B_{\mathrm{expert}}$ 是单个专家对象大小，$b_{\mathrm{cg}}$ 是 CPU-GPU 有效带宽。该值为不考虑计算与传输重叠时的专家权重传输时间估计，主要用于不同缓存策略之间的相对比较。对于包含预取的策略，模拟器额外统计预取传输量和预取浪费；正文主策略 WREC-H2 不启用预取，因此主结果中的 浪费为 $0$。

模拟器还计算 WREC 相对 LRU 的收益。设某一缓存预算下策略 $\pi$ 的传输时间估计值为 $T_\pi$，LRU 的结果为 $T_{\mathrm{LRU}}$，则策略相对 LRU 的收益为

\begin{equation}
    \label{equ:wrec-gain}
    G_{\pi}=\frac{T_{\mathrm{LRU}}-T_{\pi}}{T_{\mathrm{LRU}}}.
\end{equation}

同时，本文使用 $T_\pi-T_{\mathrm{Belady}}$ 描述在线策略与 Belady 离线上界之间的剩余差距。该指标用于判断 WREC 是否真正接近离线上界，而不仅仅是超过某一个在线基线。

\subsection{运行时接口与有限槽位原型}

除离线重放模拟器外，本文还实现了运行时影子接口（runtime shadow）。该接口将在线 WREC 策略封装为轻量旁路服务。推理框架在每个路由事件产生后向旁路服务发送请求编号、层号、词元位置和选中专家；旁路服务顺序维护 WREC 在线状态和影子驻留集合，并返回每个专家的拟准入、拟绕过和拟换出决策。该模式不控制真实专家加载，但可以验证 WREC 的状态更新和决策逻辑能否按运行时事件流在线执行。

在 vLLM 集成原型中，本文进一步实现了专家驻留管理接口。vLLM 的 PagedAttention 设计体现了通过运行时内存管理提升大语言模型服务效率的系统思路~\cite{kwon2023pagedattention}；本文原型将这一类运行时管理思想延伸到 MoE 专家权重驻留路径。该接口支持三种层次。第一是仅状态模式（state-only），只导出路由专家事件并消费旁路服务决策，不改变真实专家权重放置。第二是行复制框架模式（row-copy skeleton），为特定专家保留 CPU 备份，并在需要时把专家切片复制回 GPU 权重张量。第三是有限槽位模式（finite-slot），当 \texttt{WREC\_EXPERT\_RESIDENCY\_SLOT\_CAPACITY} 小于模型全局专家数时，将每层专家权重替换为有限 GPU 槽位张量，并维护逻辑专家到槽位编号的 \texttt{expert\_map}。MoE 内核通过该映射访问当前驻留的逻辑专家。

有限槽位原型用于验证 WREC 决策进入真实专家驻留路径的可行性，不替代第 5 章的轨迹驱动主实验。该原型要求当前批次中每层实际访问的去重专家数量不超过槽位容量，否则有限槽位无法保证正确执行。由于该路径同时受到 vLLM 调度、CPU 卸载、副本管理和内核映射约束，第 5 章仅将其作为运行时接入验证。

通过上述设计，WREC 将历史工作负载先验、请求内统计和跨层转移关系统一到在线评分函数中，并通过准入、绕过和换出机制维护有限专家驻留集合。与 MoE-Infinity、MoE-Lightning 等系统侧重专家卸载、预取和执行流水线不同，WREC 将重点放在给定路由轨迹下的在线专家缓存替换与准入决策上~\cite{xue2024moeinfinity,cao2025moelightning}。主策略采用最近历史权重为 $0$ 的 请求内/跨层 WREC-H2；离线重放模拟器用于隔离比较缓存策略，运行时影子验证和有限槽位原型用于检查接入真实推理系统的可行性。第 5 章将在此基础上给出实验结果与分析。
